\documentclass[12pt]{article}
\usepackage[utf8]{inputenc}
\usepackage[T1]{fontenc}
\usepackage[margin=1in]{geometry}
\usepackage{setspace}
\usepackage{xcolor}
\usepackage[colorlinks=true, urlcolor=blue!60!black, linkcolor=blue!60!black]{hyperref}
\setlength{\parindent}{0em}
\setlength{\parskip}{0.6em}
\renewcommand{\baselinestretch}{1.0}
\begin{document}
\begin{center}
{\LARGE \textit{Romer v. Evans}, 517 U.S. 620 (1996)\par}
\vspace{0.5em}
{\large Case Brief --- Two-Page Edition\par}
\end{center}

\section{Memory Jogger}
Colorado's Amendment 2, barring all state and local protections against discrimination for gay, lesbian, and bisexual people, failed even rational-basis review because its sweeping disability could be explained only by animus (\href{https://www.courtlistener.com/opinion/118027/romer-v-evans/}{517 U.S. 620}).

\section{Detailed Case Facts}
After Aspen, Boulder, Denver, and state agencies banned sexual-orientation discrimination, Colorado voters in 1992 approved Amendment 2, forbidding the State and every subdivision to grant homosexual, lesbian, or bisexual orientation ``any minority status, quota preferences, protected status or claim of discrimination.'' Richard Evans, other individuals, and the three cities sued.

\section{Procedural History}
The Colorado courts enjoined Amendment 2, the Colorado Supreme Court applying strict scrutiny on a political-participation theory. The U.S. Supreme Court affirmed on rational-basis grounds.

\section{Judicial Votes}
Affirmed, 6--3. Kennedy wrote for Stevens, O'Connor, Souter, Ginsburg, and Breyer; Scalia, joined by Rehnquist and Thomas, dissented.

\section{Holding}
A law singling out one class by a single trait and barring it from seeking protection at every level of government lacks a rational relationship to any legitimate purpose; a bare desire to harm a politically unpopular group is not a legitimate interest.

\section{Analysis of Opinions}
Kennedy rejected the ``special rights'' framing: Amendment 2 imposed a ``special disability upon those persons alone,'' denying protections most people take for granted. Its breadth was so discontinuous with the State's asserted interests in associational freedom and enforcement resources that it raised ``the inevitable inference that the disadvantage imposed is born of animosity toward the class of persons affected.'' Scalia called it a ``modest attempt'' to preserve traditional mores, argued that \textit{Bowers v. Hardwick} made it rational, and accused the Court of siding with ``the lawyer class.''

\section{Examples: Future Applications}
\textbf{Same-Side:} a state ban on all protections for unmarried cohabitants fails. \textbf{Opposite-Side:} a uniform statewide employment code preempting all local ordinances is valid. \textbf{Fence-Sitter:} a single city's charter amendment barring local sexual-orientation protections, which the Sixth Circuit upheld after \textit{Romer}.

\section{Critique}
Critics fault the silence about \textit{Bowers} and an unannounced ``rational basis with bite,'' and Evan Gerstmann called the reasoning ``sloppy.'' Defenders recast it as a pariah principle (Farber and Sherry), a bill of attainder (Amar), or minimalism (Sunstein).

\section{Key Quotations}
\begin{quote}
``A law declaring that in general it shall be more difficult for one group of citizens than for all others to seek aid from the government is itself a denial of equal protection of the laws in the most literal sense.''
\end{quote}
\noindent --- Justice Kennedy, Opinion of the Court, 517 U.S. at 633.

\end{document}
